\documentclass[a4paper,12pt]{article}
\usepackage[T2A]{fontenc}
\usepackage[utf8]{inputenc}
\usepackage[russian]{babel}
\usepackage{amsmath,amssymb,mathtools}
\usepackage{helvet}
\renewcommand{\familydefault}{\sfdefault}
\usepackage{icomma}
\usepackage[left=20mm,right=20mm,top=20mm,bottom=21mm]{geometry}
\usepackage{microtype}
\usepackage{tikz}
\usetikzlibrary{arrows.meta,positioning,shapes.geometric,calc}
\usepackage{tabularx,booktabs,longtable}
\usepackage{enumitem}
\usepackage{parskip}
\usepackage{fancyhdr}
\usepackage{setspace}
\usepackage{xcolor}
\usepackage[hidelinks]{hyperref}
\definecolor{accent}{HTML}{008080}
\definecolor{darkaccent}{HTML}{004D4D}
\definecolor{textgray}{HTML}{36393A}
\definecolor{lightgray}{HTML}{E0E8E6}
\color{textgray}
\setstretch{1.14}
\setlength{\parskip}{4pt}
\setlist[enumerate]{leftmargin=*,itemsep=9pt,topsep=6pt}
\pagestyle{fancy}
\fancyhf{}
\fancyhead[L]{\small\color{darkaccent}Комбинаторика: метод позиций}
\fancyhead[R]{\small Анастасия Павлова}
\fancyfoot[C]{\small\thepage}
\renewcommand{\headrulewidth}{0.4pt}
\renewcommand{\headrule}{\hbox to\headwidth{\color{lightgray}\leaders\hrule height \headrulewidth\hfill}}
\setlength{\headheight}{25pt}
\addtolength{\topmargin}{-10pt}
\newcommand{\mysection}[1]{\vspace{6pt}{\Large\bfseries\color{darkaccent}#1\par}\vspace{2pt}\textcolor{accent}{\rule{\linewidth}{0.6pt}}\vspace{3pt}}
\newcommand{\mysub}[1]{\vspace{8pt}{\large\bfseries\color{darkaccent}#1\par}\vspace{1pt}}
\newcommand{\idea}[1]{\noindent\textcolor{accent}{\textbf{Ключевая идея.}} #1}
\newcommand{\warning}[1]{\noindent\textbf{Обратите внимание.} #1}
\newcommand{\ex}[1]{\noindent\textbf{Пример.} #1}
\newcommand{\slot}[1]{\tikz[baseline=(n.base)]\node[draw=accent,very thick,rounded corners=2pt,minimum width=0.82cm,minimum height=0.70cm,inner sep=1pt,anchor=base] (n) {$#1$};}
\newcommand{\trainitem}[2]{\item \textbf{#1}\par #2}
\begin{document}
\thispagestyle{empty}
\begin{center}
\vspace*{22mm}
{\Large\bfseries\color{darkaccent}МАТЕМАТИКА\par}
\vspace{28mm}
{\fontsize{34}{39}\selectfont\bfseries Чек-лист\par}
\vspace{10mm}
{\fontsize{21}{27}\selectfont\bfseries\color{darkaccent}
Комбинаторика\\[3pt]
Метод позиций и правило произведения\par}
\vspace{13mm}
{\large Подготовка к ЕГЭ по профильной математике\par}
\vfill
{\large\bfseries Анастасия Павлова\par}
\vspace{8mm}
{\large 12 ноября 2026 г.\par}
\vfill
{\normalsize Лёвин Артём Александрович\par}
{\normalsize эксклюзивно для Анастасии Павловой\par}
\vspace{12mm}
\begin{tikzpicture}
\draw[color=accent!65,line width=1.0pt] (-6,0) .. controls (-3.9,0.55) and (-2.5,-0.45) .. (0,0)
 .. controls (2.5,0.45) and (3.9,-0.55) .. (6,0);
\end{tikzpicture}
\end{center}
\clearpage

\mysection{1. Что мы считаем в комбинаторике}
Комбинаторика помогает ответить на вопрос, \textbf{сколько различных вариантов} удовлетворяют условиям задачи. Вариантами могут быть числа, коды, маршруты, расписания, награды, очереди или порядки предметов.

\idea{Сначала нужно определить \textbf{позиции} (места, которые заполняем) и \textbf{кандидатов} (объекты, которые можно поставить на каждое место).}

\mysub{Правило произведения}
Если для первой позиции предусмотрено $a$ вариантов, а после каждого такого выбора для второй позиции имеется $b$ вариантов, общее число вариантов равно $a\cdot b$. Для нескольких позиций действуем аналогично:
\[
N=n_1\cdot n_2\cdot\ldots\cdot n_k.
\]
Здесь $n_i$ --- \textbf{количество допустимых кандидатов} на очередную позицию с учётом уже сделанных выборов. В клетках ниже записаны именно количества вариантов, а не выбранные цифры.

\begin{center}
$\slot{3}\;\slot{4}\quad\Longrightarrow\quad N=3\cdot4=12.$
\end{center}

\mysub{Двузначное число: когда повторы разрешены}
\ex{Сколько двузначных чисел можно записать цифрами $0,1,2,3$, если цифры могут повторяться?}

Для десятков доступно только $1,2,3$: \textbf{первым нуль быть не может}. Для единиц подходят все четыре цифры, включая нуль и уже использованную цифру:
\[
\underbrace{3}_{\text{десятки}}\cdot
\underbrace{4}_{\text{единицы}}=\boxed{12}.
\]
Например, числа $11$, $22$, $33$ здесь допустимы, поскольку запрета на повторение нет.

\mysub{А если цифры должны различаться?}
В той же задаче при условии «цифры различны» после выбора десятков остаётся лишь три допустимые цифры:
\[
\underbrace{3}_{\text{десятки}}\cdot
\underbrace{3}_{\text{единицы}}=\boxed{9}.
\]
\warning{Слова «различные числа» означают, что мы считаем разные результаты. Слова «числа с различными цифрами» дополнительно запрещают повтор цифр \emph{внутри одного числа}.}

\clearpage
\mysection{2. Нуль, повторы и количество разрядов}

\mysub{Число и числовой код --- разные объекты}
У трёхзначного числа первая цифра должна отличаться от нуля. У трёхсимвольного кода первая цифра может быть нулём, если условие не запрещает этого.

\ex{Из цифр $0,1,2,3$ составляем трёхзначные числа, допуская повторения:}
\[
3\cdot4\cdot4=\boxed{48}.
\]
Для трёхсимвольных кодов из тех же цифр при разрешённых повторениях получится $4\cdot4\cdot4=\boxed{64}$, поскольку код $003$ допустим.

\mysub{Повторы запрещены}
\ex{Сколько трёхзначных чисел без повторяющихся цифр можно получить из $0,1,2,3$?}
\[
\underbrace{3}_{\text{сотни}}\cdot
\underbrace{3}_{\text{десятки}}\cdot
\underbrace{2}_{\text{единицы}}=\boxed{18}.
\]
Для первой позиции нельзя взять нуль. После выбора ненулевой цифры для второй позиции остаются три цифры (включая нуль), для третьей --- две.

\mysub{Почему мы обычно не перебираем все комбинации}
Таблица удобна для \emph{проверки} небольшой задачи с двумя позициями: строки фиксируют выбор первой цифры, столбцы --- второй. С тремя и более позициями полный перебор быстро становится громоздким.

\idea{Количество вариантов, как правило, удобнее считать по позициям. Полный список нужен лишь тогда, когда его прямо требует задача.}

\mysection{3. Факториал и перестановки}
Если $n$ разных объектов нужно разместить по $n$ местам без повторений, число порядков равно
\[
n!=n(n-1)(n-2)\cdots2\cdot1.
\]
Например, $4!=4\cdot3\cdot2\cdot1=24$ и $5!=120$.

\ex{Пять разных школьных предметов нужно поставить в расписание пяти уроков, каждый предмет встречается один раз.} Для первого урока есть пять кандидатов, для второго --- четыре, затем три, два и один. Поэтому $5!=\boxed{120}$ расписаний.

Для слова «гипотенуза» все \textbf{10} букв различны, значит, можно получить $10!$ перестановок букв. У слова «треугольник» \textbf{11} различных букв, поэтому число перестановок равно $11!$.

\clearpage
\mysection{4. Как распознать позиции в текстовой задаче}

\mysub{Маршруты}
\ex{Из пункта $A$ в пункт $B$ ведут три дороги, а из $B$ в $C$ --- четыре. Каждый маршрут состоит из выбора дороги на первом и на втором участке.}
\[
\underbrace{3}_{A\to B}\cdot\underbrace{4}_{B\to C}=\boxed{12}.
\]
Здесь позиции --- \textbf{этапы пути}, кандидаты --- дороги на каждом этапе.

\mysub{Награды}
\ex{В турнире участвуют 32 команды. Одной команде присуждается золото, другой --- серебро.}
\[
\underbrace{32}_{\text{золото}}\cdot\underbrace{31}_{\text{серебро}}=\boxed{992}.
\]
Позиции --- \textbf{две разные награды}. Одна команда не может получить обе медали, поэтому количество кандидатов уменьшается.

Аналогично, если среди 18 человек распределяют три разные должности без совмещения, число способов равно $\boxed{18\cdot17\cdot16=4896}$.

\mysection{5. Условия, ограничивающие конкретные позиции}
\mysub{Заранее заданная цифра}
Если нужно составить пятизначные числа из $1,2,3,4,5$ без повторений так, чтобы \textbf{последней была цифра 3}, последняя позиция имеет одного кандидата. Четыре остальные цифры переставляем свободно:
\[
4\cdot3\cdot2\cdot1\cdot\underbrace{1}_{\text{последняя }3}
=\boxed{24}.
\]
Если же \textbf{первая цифра 5, а вторая 1}, обе позиции уже закреплены, поэтому остаются $3!=\boxed{6}$ вариантов.

\mysub{Кратность пяти}
Число делится на $5$, если оканчивается на $0$ или $5$. Из цифр $1,2,3,4,5,6$, используемых \textbf{по одному разу}, на конце подходит только $5$. Остальные переставляем:
\[
5!=\boxed{120}.
\]
Если среди разрешённых цифр есть нуль, окончания $0$ и $5$ рассматривают отдельно.

\clearpage
\mysection{6. Ограничения по чётным и нечётным местам}
В условии могут требовать, чтобы цифры \textbf{только на некоторых позициях} не совпадали. Позиции нумеруем слева направо: $1,2,3,\ldots$

\mysub{Пример: пять цифр, различны места 1, 3 и 5}
Считаем все пятизначные числа, у которых цифры на \textbf{нечётных местах} (первом, третьем, пятом) попарно различны; цифры на втором и четвёртом местах могут повторяться.

\begin{center}
$\slot{9}\;\slot{10}\;\slot{9}\;\slot{10}\;\slot{8}$
\end{center}
\[
N=9\cdot10\cdot9\cdot10\cdot8=\boxed{64\,800}.
\]
Первая позиция исключает нуль. Третья может содержать нуль, но обязана отличаться от первой. Пятая должна отличаться от первой и третьей. На второй и четвёртой ограничений нет.

\mysub{Пример: шесть цифр, различны места 2, 4 и 6}
Для шестизначных чисел при условии попарного различия цифр \textbf{на чётных местах} получаем:
\[
\underbrace{9}_{1}\cdot\underbrace{10}_{2}\cdot
\underbrace{10}_{3}\cdot\underbrace{9}_{4}\cdot
\underbrace{10}_{5}\cdot\underbrace{8}_{6}
=\boxed{648\,000}.
\]
Ограничения второй, четвёртой и шестой позиций связаны \emph{между собой}. Для первой, третьей и пятой позиций остаются обычные правила разрядности и повторений.

\mysub{Пример: нечётное трёхзначное число без повторений}
Из цифр $0,1,2,3,4,5$ нужно составить нечётные трёхзначные числа без повторения цифр. Удобно \textbf{сначала выбирать последнюю цифру}:
\[
\underbrace{3}_{\text{последняя: }1,3,5}\cdot
\underbrace{4}_{\text{первая: ненулевая и не равна последней}}\cdot
\underbrace{4}_{\text{средняя: оставшиеся цифры}}
=\boxed{48}.
\]
\warning{Для таких задач важно учитывать одновременно чётность, невозможность ведущего нуля и запрет повторений. Простое умножение одинаковых «остатков цифр» может дать неверный ответ.}

\clearpage
\mysection{7. Соседство предметов: метод блока}
\idea{Если несколько определённых предметов должны стоять рядом, временно «склейте» их в один блок. Сначала расставьте блок и остальные предметы, затем учтите перестановки внутри блока.}

\mysub{Два тома рядом}
\ex{На полке есть шесть отдельных книг и два различающихся тома одного произведения. Оба тома обязательно должны стоять рядом, но их взаимный порядок произволен.}

Пара томов становится одним блоком: вместе с шестью отдельными книгами имеется \textbf{семь объектов}. Внутри блока два тома можно переставить двумя способами:
\[
\underbrace{7!}_{\text{порядок семи объектов}}
\cdot\underbrace{2!}_{\text{порядок двух томов}}
=5040\cdot2=\boxed{10\,080}.
\]

\mysub{Три тома рядом}
\ex{На той же полке восемь книг: пять отдельных книг и три различающихся тома одной серии, которые должны стоять подряд.} Получаем шесть объектов (блок + пять книг), а внутри блока можно переставлять три тома:
\[
6!\cdot3!=720\cdot6=\boxed{4320}.
\]

\mysub{Краткая памятка перед решением}
\begin{enumerate}[label=$\square$,leftmargin=1.4cm,itemsep=4pt]
\item Я определил(а), что именно считаю: число, код, маршрут, расписание или порядок предметов.
\item Я выбрал(а) позиции и понял(а), какие объекты их заполняют.
\item Я проверил(а), допускаются ли повторы и можно ли поставить нуль первым.
\item Я учёл(а) фиксированные места, чётность, кратность и условие соседства.
\item Я проверил(а), не посчитал(а) ли один и тот же вариант дважды.
\end{enumerate}

\mysub{Типичные ошибки}
\begin{itemize}[leftmargin=*,itemsep=1pt,topsep=2pt]
\item Принимать число кандидатов на позиции за конкретную цифру.
\item Разрешать повторения там, где условие требует разных цифр.
\item Запрещать начальный нуль даже для кода, где он допустим.
\item Забывать, что в задаче с делимостью последнюю цифру полезно выбирать первой.
\item «Склеивать» тома, забывая перемножить на число перестановок внутри блока.
\end{itemize}

\clearpage
\thispagestyle{fancy}
\begin{center}
{\Large\bfseries\color{darkaccent}Алгоритм: метод позиций}\par
\vspace{9mm}
\begin{tikzpicture}[
  >=Latex, line width=.85pt,
  every node/.style={align=center,font=\normalsize},
  start/.style={ellipse,draw=accent,line width=1pt,minimum width=5.0cm,minimum height=0.95cm},
  action/.style={rectangle,rounded corners=3pt,draw=accent,line width=.85pt,text width=9.4cm,minimum height=1.0cm,inner sep=7pt},
  decision/.style={diamond,aspect=3.0,draw=accent,line width=.85pt,text width=5.0cm,inner sep=4pt},
  arrow/.style={->,>=Latex,line width=1pt,draw=darkaccent}]
\node[start] (start) at (0,0) {Начало: внимательно прочитайте условие};
\node[action] (a) at (0,-2.0) {Определите, \textbf{что} считаем, и нарисуйте позиции: разряды числа, медали, места, этапы пути};
\node[decision] (d) at (0,-4.55) {Есть ограничения на позиции?};
\node[action,text width=5.9cm,minimum height=1.7cm] (yes) at (4.65,-7.0) {Уточните: нуль первым, повторы, чётность, фиксированные места, соседство};
\node[action] (b) at (0,-9.35) {Для каждой позиции посчитайте \textbf{допустимых кандидатов} с учётом предыдущих выборов};
\node[action] (c) at (0,-11.65) {Перемножьте количества: $N=n_1\cdot n_2\cdots n_k$};
\node[action] (check) at (0,-13.95) {Проверьте: учтены ли все условия, нет ли повторного подсчёта, разумен ли результат?};
\node[start] (end) at (0,-16.0) {Ответ};
\draw[arrow] (start)--(a);
\draw[arrow] (a)--(d);
\draw[arrow] (d.east) -| node[pos=.25,above,font=\small]{да} (yes.north);
\draw[arrow] (yes.south) |- (b.east);
\draw[arrow] (d.south) -- node[right,font=\small]{нет} (b.north);
\draw[arrow] (b)--(c);
\draw[arrow] (c)--(check);
\draw[arrow] (check)--(end);
\end{tikzpicture}
\end{center}
\clearpage

\mysection{Тренировочный блок}
\textbf{Путь А. 10 задач.} Условия сформулированы полностью. Во всех задачах нужно определить позиции, указать число вариантов для каждого выбора и вычислить итог. Условия читаются буквально: разрешение повторений, возможность начального нуля и порядок объектов имеют значение.

\mysub{Базовый уровень}
\begin{enumerate}[label=\textbf{\arabic*.},leftmargin=9mm,itemsep=12pt]
\item \textbf{Код электронного шкафчика.} В спортивном центре каждый шкафчик открывается двухсимвольным цифровым кодом. Для каждой позиции разрешены только цифры $0$, $2$, $4$ и $7$. Одну и ту же цифру можно использовать несколько раз; код, начинающийся с нуля, действителен. Сколько разных кодов можно составить?

\item \textbf{Номера экспонатов.} Организаторы выставки хотят присваивать экспонатам двузначные номера, используя только цифры $0$, $2$, $4$ и $7$. Номер должен быть именно двузначным числом: запись $02$ недопустима. Повторять цифры разрешено. Сколько разных номеров подходит под эти правила?

\item \textbf{Номера с разными цифрами.} Для другой части выставки организаторы используют те же цифры $0$, $2$, $4$, $7$, но теперь запрещают повторять одну и ту же цифру внутри номера. Номера по-прежнему должны быть двузначными числами. Сколько таких номеров существует?

\item \textbf{Пешеходный маршрут.} Турист отправляется из посёлка $A$ к смотровой площадке $B$, а затем идёт к маяку $C$. Между $A$ и $B$ есть четыре разные тропы, между $B$ и $C$ --- три. Турист выбирает ровно одну тропу для каждого из двух участков, и любые такие сочетания разрешены. Сколько разных маршрутов из $A$ в $C$ через $B$ возможно?

\item \textbf{Трёхзначные числа.} Для математической игры нужно составлять трёхзначные числа из цифр $0$, $1$, $3$ и $5$. В записи каждого числа цифры не должны повторяться; нуль в разряде сотен запрещён. Использовать все четыре цифры одновременно не требуется. Сколько различных трёхзначных чисел можно получить?
\end{enumerate}

\clearpage
\mysection{Тренировочный блок — продолжение}
\mysub{Уровень выше среднего}
\begin{enumerate}[label=\textbf{\arabic*.},start=6,leftmargin=9mm,itemsep=13pt]
\item \textbf{Число, кратное пяти.} На конкурсе участнику выдают шесть карточек с цифрами $1$, $2$, $3$, $4$, $5$ и $6$. Требуется разложить все шесть карточек в ряд, чтобы получилось шестизначное число, делящееся на $5$. Каждая карточка используется ровно один раз. Сколько различных шестизначных чисел можно составить?

\item \textbf{Расписание занятий.} Школьный класс в один день изучает пять разных предметов: математику, русский язык, физику, историю и литературу. По каждому предмету должен пройти ровно один урок. Все пять уроков стоят подряд, и ограничений на их порядок нет. Сколько различных расписаний можно составить на этот день?

\item \textbf{Награждение участников.} В финал школьной олимпиады вышли восемь человек. Организаторы должны вручить одну золотую, одну серебряную и одну бронзовую медаль. Каждый призёр получает ровно одну медаль, а порядок наград важен: золото и серебро нельзя поменять местами, сохранив тот же результат. Сколькими способами можно распределить три медали?

\item \textbf{Различие цифр на чётных местах.} Рассматривают все четырёхзначные числа в десятичной записи. Известно только одно дополнительное условие: цифры, стоящие \textbf{на втором и четвёртом местах}, должны отличаться друг от друга. На первой и третьей позициях цифры могут совпадать с любыми другими, если число остаётся четырёхзначным. Сколько чисел удовлетворяют условию?
\end{enumerate}

\mysub{Сложный уровень}
\begin{enumerate}[label=\textbf{\arabic*.},start=10,leftmargin=9mm,itemsep=13pt]
\item \textbf{Трёхтомник на книжной полке.} В библиотеке нужно расставить в один ряд девять различных книг. Шесть из них являются самостоятельными произведениями разных авторов, а оставшиеся три --- первый, второй и третий тома одной серии. По просьбе библиотекаря все три тома должны стоять \textbf{рядом}, но внутри этой тройки порядок томов может быть любым. Сколько различных расстановок всех девяти книг существует?
\end{enumerate}

\clearpage
\mysection{Ответы к тренировочному блоку}
\textit{В таблице указаны выражения для подсчёта и окончательные ответы. Полные решения выполняйте самостоятельно по методу позиций.}
\vspace{5mm}
\renewcommand{\arraystretch}{1.38}
\begin{center}
\begin{tabularx}{\linewidth}{@{}c X r@{}}
\toprule
\textbf{№} & \textbf{Подсчёт} & \textbf{Ответ} \\
\midrule
1 & $4\cdot4$ & $16$ \\
2 & $3\cdot4$ & $12$ \\
3 & $3\cdot3$ & $9$ \\
4 & $4\cdot3$ & $12$ \\
5 & $3\cdot3\cdot2$ & $18$ \\
6 & $5!$ & $120$ \\
7 & $5!$ & $120$ \\
8 & $8\cdot7\cdot6$ & $336$ \\
9 & $9\cdot10\cdot10\cdot9$ & $8100$ \\
10 & $7!\cdot3!$ & $30\,240$ \\
\bottomrule
\end{tabularx}
\end{center}
\vspace{11mm}
\mysection{Итоговая памятка}
\renewcommand{\arraystretch}{1.5}
\begin{tabularx}{\linewidth}{@{}p{54mm} X@{}}
\toprule
\textbf{Сигнал в условии} & \textbf{Что проверить или сделать}\\
\midrule
Число и цифра $0$ & Первый разряд ненулевой.\\
Код и цифра $0$ & Начальный нуль обычно допустим.\\
«Цифры различны» & После выбора кандидатов становится меньше.\\
«Каждый объект один раз» & При расстановке всех разных объектов появляется факториал.\\
Чётность или кратность & Удобно начать с последней цифры и проверить исключения.\\
Различны цифры на отдельных местах & Запрет относится только к указанным позициям.\\
Предметы стоят рядом & Объединить их в блок и учесть перестановки внутри.\\
\bottomrule
\end{tabularx}
\vfill
\begin{center}
\small\color{darkaccent}Лёвин Артём Александрович эксклюзивно для Анастасии Павловой
\end{center}
\end{document}
